\section{Introduction}
Causal inference has emerged as a central paradigm for understanding and influencing complex real-world systems in domains such as medicine \citep{Shi2022LearningCEA,Smith2023ApplicationOTA,Cai2017IdentificationOAA,Prosperi2020CausalIAA}, natural sciences \citep{Varghese2024CausalDTA,Zucker2021LeveragingSBA,Jiao2018BivariateCDA}, engineering \citep{Chadbourne2023ApplicationsOCA,Graham2022CausalIFA,Ji2022PerfcePDA,Li2024InferringHTA,Siebert2022ApplicationsOSA}, economics \citep{Abbring2025PhilipGWA,Demeter2025AssessingTIA,Athey2016TheSOA}, and education \citep{Nagai2024EducationalEIA,Chi2022APGA,weissman2021methods}. 
Over the past decade, substantial progress has been made in \emph{static} causal estimation—learning causal effects or graph structures from a \emph{fixed}, passively observed dataset \citep{johansson2016learning,shalit2017estimating,alaa2017bayesian,alaa2018limits,yoon2018ganite,hassanpour2019counterfactual,zhang2020learning,hassanpour2019learning,assaad2021counterfactual,curth2021nonparametric,zheng2018dags,yu2019dag,annadani2023bayesdag,cundy2021bcd,kalainathan2022structural,lippe2021efficient,huang2020causal,shimizu2006linear,hoyer2008nonlinear,chickering2002optimal,chickering2015selective,zhu2019causal,buhlmann2014cam,spirtes2000causation,kaptanoglu2021pysindy,Brunton2015DiscoveringGEA,Shojaee2023TransformerbasedPFA,Valipour2021SymbolicGPTAGA,Petersen2019DeepSRA,Tegmark2019AIFAA,Udrescu2020AIF2A,Cranmer2020DiscoveringSMA,Biggio2021NeuralSRA}. 
However, in many realistic settings, reliable inference is only possible through \emph{active intervention and experimentation}—the agent must decide \emph{what to observe next} rather than consume data passively \citep{von2019optimal,Seitzer2021CausalIDA,Mehrjou2020CausalCRA}. 
This introduces a fundamentally different challenge: causal inference becomes a \emph{sequential decision process under resource constraints}, where experimental efficiency matters as much as correctness.

A crucial activity in causal inference as a discipline is the active design of interventions \citep{von2019optimal,Seitzer2021CausalIDA,Mehrjou2020CausalCRA}.
Without some exceptions, it is not generally possible to conclude on causal relationships without systematic interventions \citep{pearl2009causality}, e.g., in the form of treatments or systematic exploration of an environment.
Since the execution of such experiments is often very costly in practice, properly designing and executing them is arguably just as important as the final inference activity.
In fact, without sufficient information from well-design experiments, no inference might be possible.

Despite its importance, this interactive dimension is almost entirely absent \me{We are aware of} from existing benchmarks. 
Frameworks for causal graph discovery \citep{chengcausaltime,lawrence2021data,runge2019causality,chevalley2023causalbench}, symbolic regression \citep{olson2017pmlb,la2021contemporary,romano2021pmlb}, and treatment effect estimation \citep{CauseBox2021,ACCBench2024,CIPCaDBench2022,CausalDiscoveryFramework2024,RealCause2021,BayesianCausalDiscovery2023,CausalBench2024,EvaluationFramework2022,IBM2022,DoWhy2021,Framework2021} almost exclusively evaluate models on a fixed dataset, scoring only the final outcome rather than \emph{how the knowledge was acquired}. 
The only known interactive benchmark, CausalWorld \citep{ahmed2020causalworld}, focuses on robotic control and does not \emph{explicitly} evaluate the correctness of causal reasoning itself. 
Thus, there is currently no rigorous framework for benchmarking \emph{agent-based causal inference} as a problem of strategic experimental design.

To close this gap, we introduce \game, a turn-based, agent-centric causal inference framework in which players actively select interventions, gather observations, and decide when to stop and return a causal conclusion. 
As argued by \cite{BayesianCausalDiscovery2023}, the distinction between causal graph discovery and treatment effect estimation is artificial, in particular since advanced treatment effect estimation techniques might solve causal graph discovery as a sub-task.
To \game, these are simply different \emph{missions} with respective \emph{result metrics}, e.g., F1 or SHD for graph discovery, and PEHE for treatment effect estimation.
Unlike static benchmarks, \game jointly evaluates (i) the \emph{quality of causal conclusions} and (ii) the \emph{efficiency of the experimental strategy}. 
The accompanying Python package \package provides a rigorously modular implementation, enabling pluggable agents, structural causal models (SCMs), behavior and result metrics, stopping functions, and fully seedable execution for reproducible comparison across agents and settings.

\textbf{Contributions.} This work introduces:
\begin{itemize}
    \item A formalized \emph{agent-based causal inference game} that evaluates not only \emph{what} an agent concludes, but \emph{how intelligently} it gathers information.
    \item A modular, extensible \emph{system architecture} implemented in \package, supporting dynamic SCM generation from datasets, Bayesian networks, and synthetic rules.
    \item A principled framework for multi-mission evaluation, covering DAG recovery, CATE estimation, SCM reconstruction, and beyond.
    \item A fully seedable execution model enabling \emph{fair, reproducible, and comparable} evaluation of heterogeneous causal agents.
\end{itemize}

By reframing causal inference as a strategic experimental process rather than a passive estimation task, \game enables evaluation of the next generation of \emph{interactive causal agents}—those that do not merely estimate causality, but \emph{actively learn it}.
